% ============================================================
%  SECTION 5: EMPIRICAL STRATEGY
% ============================================================

\section{Empirical Strategy}
\label{sec:empirical_strategy}

\subsection{Primary Estimator: Callaway and Sant'Anna (2021)}

Our primary estimator is the group-time average treatment effect on the
treated (ATT) proposed by \textcite{callaway2021difference}. This
estimator avoids the negative weighting and heterogeneity bias that
can afflict two-way fixed effects specifications when treatment effects
vary over time \parencite{goodmanbacon2021difference,
sun2021estimating}.

In our setting, treatment timing is sharp: all Tennessee public colleges
are treated simultaneously at $t^* = 1986$ (the first birth cohort with
unambiguous HOPE exposure). The Callaway--Sant'Anna estimator computes
group-time ATTs by comparing the outcome evolution of the treated group
to a ``never-treated'' comparison group. We aggregate the group-time
ATTs into an overall ATT and dynamic event-study coefficients.

We report TWFE event-study estimates for comparability with the prior
difference-in-differences literature
\parencite{ashenfelter1985panel, card1994minimum}. The TWFE
specification is:

\begin{equation}
y_{ct} = \alpha_c + \lambda_t + \sum_{k=-5}^{5} \beta_k \cdot
\mathbf{1}[t - t^* = k] \cdot \text{Treated}_c + \varepsilon_{ct}
\label{eq:twfe}
\end{equation}

\noindent where $y_{ct}$ is the outcome for college $c$ in birth cohort
$t$, $\alpha_c$ are college fixed effects, $\lambda_t$ are cohort fixed
effects, $\text{Treated}_c$ indicates Tennessee public colleges, and
$k = -1$ is the omitted reference period.

\subsection{Identification}

Identification requires that Tennessee public colleges would have
followed parallel trends in socioeconomic composition to control
colleges absent the HOPE Scholarship. We assess this assumption in
three ways.

First, we inspect pre-treatment event-study coefficients. Under the
national control group, coefficients at $k = -5$ through $k = -3$ show
some divergence from zero for bottom-quintile share, raising concerns about
longer-horizon pre-trends. The $k = -5$ coefficient is driven by a
single cohort (1980 births) and its magnitude is sensitive to the
inclusion of this endpoint.

Second, we implement a bordering-state control group that restricts
comparisons to colleges in states adjacent to Tennessee. This
specification passes pre-trend tests: pre-treatment coefficients are
small, statistically insignificant, and show no systematic pattern.
The bordering-state specification is our preferred identification
strategy.

Third, we implement a conservative control group that drops states with
any policy changes in higher education funding during the analysis
window. Results are quantitatively similar to the baseline.

\subsection{Inference}

Standard errors are clustered at the state level, the level at which
treatment is assigned. With 27 clusters (26 control states plus
Tennessee), we are near the boundary where cluster-robust inference
becomes unreliable. In the Callaway--Sant'Anna framework, we use
the analytical standard errors provided by the \texttt{csdid} package,
which account for the two-step estimation.

For the bordering-state specification, which has only four state
clusters ($G_1 = 1$ treated), the wild cluster bootstrap is
uninformative \parencite{mackinnon2017wild}. We therefore supplement
analytical standard errors with a Fisher-style permutation test that
assigns placebo treatment to each state in turn and ranks Tennessee's
estimate against the resulting distribution. We report permutation
$p$-values in Section~\ref{sec:robustness_permutation}.

\subsection{Sensitivity Analysis}

We implement the honest confidence intervals of
\textcite{rambachan2023more}, which relax the parallel trends assumption
by allowing for bounded violations. We report results under two restrictions:

\begin{itemize}
\item \textbf{Relative magnitude} ($\Delta_\text{RM}$): post-treatment
  violations of parallel trends are bounded by a multiple $\bar{M}$ of
  the maximum pre-treatment violation. Results remain significant at
  $\bar{M} \leq 2$ for bottom-quintile share.

\item \textbf{Smoothness} ($\Delta_\text{SD}$): consecutive violations
  of parallel trends are bounded in their rate of change. Under this
  restriction, confidence intervals include zero at moderate bounds,
  indicating fragility. We interpret this as a limitation: if
  differential trends can accelerate smoothly, our estimates cannot be
  distinguished from pre-existing dynamics.
\end{itemize}

We discuss these results in detail in Section~\ref{sec:robustness} and
interpret the tension between $\Delta_\text{RM}$ and $\Delta_\text{SD}$
as informative about the nature of potential confounds.

\subsection{Synthetic Control}

As a complementary identification strategy, we implement the synthetic
control method \parencite{abadie2010synthetic}, which constructs a
data-driven counterfactual for Tennessee by selecting weighted
combinations of donor states that best match Tennessee's pre-treatment
composition trends. The donor pool consists of 22 non-merit states,
matching on pre-treatment bottom-quintile share levels (birth cohorts
1980--1984). This approach avoids imposing equal weights across all
control states and directly addresses concerns about the suitability of
any particular comparison group. We report the synthetic control trends
in Section~\ref{sec:results} and placebo-in-space inference in
Section~\ref{sec:robustness_synth}.
